\chapter{Eczema Treatment}
\index{Eczema Treatment}

\section*{Definition and Overview}
Eczema treatment aims to control itching, heal the skin, reduce flare-ups, and maintain healthy skin between episodes. Treatment is stepped, beginning with daily moisturising and good skincare, adding anti-inflammatory creams during flare-ups, and progressing to stronger treatments or specialist therapies for severe eczema. This chapter explains the full range of eczema treatments, how they work, and how to use them correctly, because proper use of simple treatments is the key to controlling the condition.

\section*{Epidemiology}
Eczema affects a large proportion of children and many adults, and under-treatment is common: many families either avoid effective treatments through fear of side effects or use treatments incorrectly. Well-managed eczema is the exception rather than the rule, and this contributes to needless suffering, sleep loss, infection, and school absence. Understanding the treatment ladder helps people use the right treatment at the right time.

\section*{Aetiology and Pathophysiology}
Treatment works by targeting the two core problems of eczema: the dry, leaky skin barrier and the inflamed immune response.
\begin{itemize}
    \item \textbf{The skin barrier:} moisturisers and emollients repair and protect the barrier, reducing dryness and irritation
    \item \textbf{Inflammation:} topical corticosteroids and other anti-inflammatory agents dampen the immune response in the skin
    \item \textbf{Itch:} controlling inflammation and barrier function breaks the itch-scratch cycle
    \item \textbf{Infection:} treating bacterial infection of broken skin is part of flare management
    \item \textbf{Severe disease:} eczema that does not respond to topical treatment involves systemic inflammation that responds to oral or injected therapies
\end{itemize}

\section*{Clinical Features}
The treatment needs reflect the severity of the eczema.
\begin{itemize}
    \item Mild eczema: small areas of dry, itchy skin, managed with moisturisers and occasional mild steroid
    \item Moderate eczema: larger or more inflamed areas requiring regular topical steroids
    \item Severe eczema: widespread, persistent, or constantly relapsing eczema requiring specialist treatment
    \item Signs that treatment is needed: persistent itch, sleep disturbance, weeping or crusted skin
    \item Signs of undertreatment: thick, lichenified skin and recurrent infection
\end{itemize}

\section*{Differential Diagnosis}
Before treating, the diagnosis should be confirmed and other conditions considered.
\begin{itemize}
    \item Contact dermatitis, requiring identification of the specific trigger
    \item Seborrhoeic dermatitis, treated differently
    \item Psoriasis, which responds to different treatments
    \item Fungal infections, which require antifungal treatment
    \item Scabies, requiring specific anti-mite treatment
\end{itemize}

\section*{When to Seek Medical Care}
People with eczema should see a doctor for diagnosis, for a treatment plan, when flare-ups are severe or not responding, and if infection is suspected. Referral to a dermatologist is appropriate for severe or treatment-resistant eczema.

\textbf{Contact your local emergency services for:}
\begin{itemize}
    \item Sudden severe spreading redness with fever
    \item Blisters with fever, which may indicate eczema herpeticum
    \item Signs of a severe allergic reaction
    \item A very unwell child with poor feeding or lethargy
\end{itemize}

\section*{Diagnosis and Investigations}
Treatment planning begins with assessment of severity and triggers.
\begin{itemize}
    \item \textbf{History:} pattern of the rash, response to previous treatment, and impact on sleep and life
    \item \textbf{Examination:} extent and severity of the eczema, and signs of infection
    \item \textbf{Severity scoring:} structured tools guide treatment decisions
    \item \textbf{Swabs:} for suspected infection
    \item \textbf{Allergy assessment:} when triggers are unclear
\end{itemize}

\section*{Management}
Treatment is stepped, and most people need a combination of approaches.
\begin{itemize}
    \item \textbf{Emollients:} the foundation; applied liberally and frequently, using creams, ointments, or bath additives
    \item \textbf{Topical corticosteroids:} used in flare-ups; potency is chosen for the body area and severity, and used for short courses or with a maintenance schedule
    \item \textbf{Topical calcineurin inhibitors:} non-steroid options for the face, neck, and skin folds
    \item \textbf{Antihistamines:} for itching and sleep
    \item \textbf{Treating infection:} antibiotic or antiseptic treatment for infected eczema
    \item \textbf{Wet wraps:} for severe flare-ups, under medical supervision
    \item \textbf{Phototherapy:} ultraviolet treatment for moderate to severe eczema
    \item \textbf{Systemic medicines:} immunosuppressants such as methotrexate or ciclosporin
    \item \textbf{Biologic therapy:} targeted injections such as dupilumab for severe eczema
\end{itemize}

\section*{Complications}
\begin{itemize}
    \item Steroid phobia leading to undertreatment and unnecessary suffering
    \item Skin thinning with prolonged misuse of strong steroids, which is uncommon with correct use
    \item Recurrent infection from untreated breaks in the skin
    \item Side effects of systemic treatments, which require monitoring
    \item Persistent sleep loss and effects on schooling and quality of life
\end{itemize}

\section*{Prognosis and Outcomes}
Eczema treatment is highly effective when used correctly. Most people achieve excellent control with moisturisers and topical steroids, and severe eczema now responds well to modern systemic and biologic treatments. The key to good outcomes is a written treatment plan, correct application, early treatment of flare-ups, and regular review.

\section*{Prevention and Self-Care}
\begin{itemize}
    \item Moisturise daily, even when the skin looks clear
    \item Apply steroid creams to inflamed skin only, and stop or step down when it settles
    \item Apply moisturiser after topical treatment, or at a different time, as advised
    \item Use the "fingertip unit" measure to apply the right amount of steroid
    \item Avoid triggers such as soaps, wool, and overheating
    \item Keep a treatment plan and review it with your doctor
    \item Seek specialist review if eczema is not improving
\end{itemize}

\section*{Sources and Further Reading}
The following reputable sources provide further information for healthcare students and patients seeking reliable, current advice about eczema treatment.
\begin{itemize}
    \item Eczema Association of Australasia. \textit{Treatment information and support.} https://www.eczema.org.au/
    \item Healthdirect Australia. \textit{Eczema treatment.} https://www.healthdirect.gov.au/eczema
    \item Australasian Society of Clinical Immunology and Allergy. \textit{Eczema treatment advice.} https://www.allergy.org.au/
    \item NHS. \textit{Atopic eczema: treatment.} https://www.nhs.uk/conditions/atopic-eczema/
    \item Mayo Clinic. \textit{Atopic dermatitis: treatment and drugs.} https://www.mayoclinic.org/diseases-conditions/atopic-dermatitis-eczema/
    \item MSD Manual. \textit{Atopic dermatitis: treatment.} https://www.msdmanuals.com/
\end{itemize}
